\begin{statement}{C.32}{Proved}{The full metabelian arithmetic-normal quotient}
Let
\[
 \rho_M:\widehat{\mathrm{GT}}_0\longrightarrow\Aut(M)
\]
be the homomorphism obtained by taking the actual normalized
representatives of B.7, and put
$G_M=\rho_M(\widehat{\mathrm{GT}}_0)$.
For either $H=G_{\mathbb Q}$ or
$H=\widehat{\mathrm{GT}}$, let $N_{M,H}$ be the closed normal
subgroup of $G_M$ generated by $\rho_M(H)$.
Then the entire quotient, without further abelianization or
reduction modulo two, is noncanonically isomorphic to
\[
                G_M/N_{M,H}\simeq\prod_{n\geq1}C_2
                \tag{C.32.1}
\]
as a topological group.
In particular every finite continuous quotient of this group is
an elementary abelian two-group. Equivalently, every continuous
homomorphism from $G_M$ to a finite group which kills $\rho_M(H)$
has elementary abelian two-group image.
This statement concerns the actual metabelian image of the full
gentle group. It does not determine the image of $\ker(\rho_M)$ in the
arithmetic-normal quotient of the full gentle group.
\end{statement}
\begin{proof}
The normalization in B.7 is a group section, not merely a choice
of representatives as a set. Thus $\rho_M$ is a continuous
homomorphism into actual automorphisms, and its compact image
$G_M$ is closed. The multiplier of an element of $G_M$ is read
from its scalar action on
$M^{\mathrm{ab}}=\widehat{\mathbb Z}^{\,2}$.
Write $U_M$ for the kernel of this multiplier.

We first check both the group law and its faithfulness on $U_M$.
For a normalized element of multiplier one, B.7 writes
\[
 x\longmapsto x,\qquad
 y\longmapsto(c_0^r)^{-1}y c_0^r,\qquad
 B=1+(X-1)(Y-1)r .
 \tag{C.32.2}
\]
The product $(X-1)(Y-1)$ is a non-zero-divisor in
$\Lambda=\widehat{\mathbb Z}[[\widehat{\mathbb Z}^{\,2}]]$,
as proved in B.7. Hence $B$ uniquely determines $r$ and the
actual normalized automorphism. The all-multiplier crossed law
(C.30.5) reduces, on $U_M$, to ordinary multiplication
$B_{uv}=B_uB_v$. Consequently
\[
                   U_M\hookrightarrow\Lambda^\times,\qquad
                   u\longmapsto B_u
                   \tag{C.32.3}
\]
is an injective homomorphism. In particular $U_M$ is abelian.
We require only the image of full gentle automorphisms here;
we do not replace it by all abstract peripheral units of $M$.

Let $c=(-1,1)$ be the specified complex conjugation. It belongs
to the Galois image and therefore to either subgroup $H$ in the
statement. Its normalized action is
$x\mapsto x^{-1}$, $y\mapsto y^{-1}$.
The entire coefficient of its action on $c_0$ is
$S_{-1}(X)S_{-1}(Y)=X^{-1}Y^{-1}$, whereas its normalized
coefficient is
\[
                         B_c=1.
\]
Since $c^2=1$, applying (C.30.5) twice gives
\[
 B_{\rho_M(c)u\rho_M(c)^{-1}}(X,Y)
                    =B_u(X^{-1},Y^{-1})=B_u^*.
 \tag{C.32.4}
\]
Every $u\in U_M$ comes from a multiplier-one element of the
full gentle group: the multiplier of any preimage is already
determined by its image on $M^{\mathrm{ab}}$.
The necessary part of B.17 therefore applies and says
$B_uB_u^*=1$. By the faithful identification (C.32.3),
(C.32.4) is the equality of actual normalized automorphisms
\[
               \rho_M(c)u\rho_M(c)^{-1}=u^{-1}
                         \qquad(u\in U_M).
 \tag{C.32.5}
\]
This avoids any inference from a nonfaithful commutator
coefficient or from an unspecified inner normalization.

Put $Q=G_M/N_{M,H}$. The cyclotomic multiplier on $H$ is onto
$\widehat{\mathbb Z}^{\times}$, because it is already onto
on $G_{\mathbb Q}$. Given $g\in G_M$, choose $h\in\rho_M(H)$
with the same multiplier. Then $gh^{-1}\in U_M$ and $g$ and
$gh^{-1}$ have the same image in $Q$. Thus the continuous
map $U_M\to Q$ is onto. It follows from the abelianity of
$U_M$ that $Q$ is abelian. Moreover $\rho_M(c)$ maps to
the identity in $Q$. Equation (C.32.5) now gives
$\bar u=\bar u^{-1}$ for every $\bar u\in Q$.
The whole compact quotient $Q$ therefore has exponent two.

It remains to determine its topological rank. Each character
$\vartheta_p$ of C.30 depends only on the normalized coefficient
$B_g$, hence factors continuously through $\rho_M$.
For continuity of this factorization one may use that a
continuous surjection from a compact space to a Hausdorff space
is a quotient map. The characters vanish on $H$, so they descend
to $Q$. The actual full witnesses $u_p$ of C.30 retain
$\vartheta_q(u_p)=\delta_{pq}$ after passage to $G_M$.
They give infinitely many independent continuous characters
of $Q$.

On the other hand $G_M$, and therefore $Q$, is a quotient of
the compact second-countable group
$\widehat{\mathrm{GT}}_0$. The countability argument in C.30
shows that $\Hom_{\mathrm{cts}}(Q,\mathbb F_2)$ has at most
countable dimension. It consequently has dimension exactly
$\aleph_0$. A profinite abelian group of exponent two is
separated by its continuous characters to $\mathbb F_2$:
each nontrivial element survives in a finite quotient, and
that quotient is a finite-dimensional $\mathbb F_2$-vector
space. Choose a countable basis of the character space.
Its evaluation map embeds $Q$ in $\prod_{n\geq1}C_2$.
Every finite projection is onto, by linear independence;
compactness then makes the whole evaluation map onto.
This proves (C.32.1).

The two choices of $H$ give the same abstract topological
isomorphism type. No equality of their closed normal subgroups,
or isomorphism induced by the natural quotient map between
them, is asserted.
\end{proof}
